\section{Independent Transition Validation for Direct Compilation}
\label{sec:direct-trace-validation}

Transition-level differential testing against the production interpreter checks a stronger observable than final-state equivalence, but it still treats one production implementation as the reference oracle. Beta~14 adds a third, deliberately smaller execution path that consumes the compiler's Change IR and independently reconstructs the expected transition sequence for the currently supported direct WebAssembly subset.

The validator first assigns deterministic lexical identifiers to structured \texttt{CHANGE} sites. A site descriptor records its scope, source line when available, target, and source operation descriptors. Site identifiers are generated by the validator rather than accepted from the WebAssembly lowerer. The validator then executes the supported IR fragment with its own expression evaluator and control-flow/call semantics, producing an expected sequence
\[
T_{IR}=\langle\langle s_i,x_i,b_i,a_i\rangle\rangle_i,
\]
where $s_i$ is the independently reconstructed Change-site identifier, $x_i$ the persistent target, and $b_i,a_i$ the numeric before/after values.

Direct WebAssembly execution independently produces the observed host trace
\[
T_{Wasm}=\langle\langle x_i,b_i,a_i\rangle\rangle_i.
\]
The translation-validation gate checks equal length and ordered equality of target, before, and after values. On success, observed occurrences can be annotated with the validator's independently reconstructed site identities. The validator does not invoke the production interpreter and does not reuse the direct backend's expression compiler.

The supported validator semantics currently cover top-level numeric creation, numeric \texttt{set/add/remove/clear}, structured conditionals, literal repetition with the 1-based \texttt{count} local, and acyclic numeric recipe calls with ranged parameters. Its expression evaluator separately implements the numeric and boolean subset used by direct compilation. CI additionally performs tamper tests that remove, reorder, or modify observed transitions and requires the validator to reject them.

This arrangement gives the direct backend two dynamic semantic checks:
\[
T_{Wasm}=T_{Interpreter}
\]
and
\[
\pi(T_{IR})=T_{Wasm},
\]
where $\pi$ erases the validator-only Change-site identifier. The second relation is particularly useful because its expected transitions are derived from Change IR rather than production interpreter history.

We describe this as \emph{translation validation evidence}, not verified compilation. The validator trusts the supplied Change IR and therefore does not prove correctness of source parsing or source-to-IR lowering. It also validates concrete executions rather than statically proving the generated WebAssembly instruction sequence correct for all inputs. The host trace callback remains part of the trusted runtime boundary.

Nevertheless, this validation boundary aligns closely with Patch's research hypothesis. State-Change Factorization makes semantic transitions explicit in the IR, while the direct runtime exposes the same before/after transition boundary. This allows backend validation to operate on semantic changes rather than only on byte-level code generation or final values.

The next strengthening is effect-aware validation. The IR-side site contract already contains operation descriptors, whereas the observed runtime event currently contains only target and before/after state. Independently checking the expected semantic effect identity, and subsequently connecting this relation to the Lean \texttt{SourceExecutes}/\texttt{Executes} models, would narrow the gap between the mechanized core and the executable backend without requiring immediate verification of the entire production compiler.
